\documentclass{article}
\usepackage{amsmath,amssymb}
\usepackage[margin=1in]{geometry}
\usepackage{mdframed}
\usepackage{needspace}

\begin{document}


\begin{flushright}
Neil Hwang\\
Linear Algebra\\
May 13, 2026\\
Homework 2\\

\end{flushright}




\bigskip
\begin{center}
\rule{0.5\textwidth}{1pt}\\[6pt]
\textbf{Homework from Lecture on May 13, 2026}\\[2pt]
\rule{0.5\textwidth}{1pt}
\end{center}
\bigskip

%% ============================================================
%% Problem 1: Change of basis verification
%% ============================================================

\needspace{4\baselineskip}
\hfill \break  $~\!\!$ \dotfill \medskip \\
\textbf{Question 1.}\\
Check by example the theorem:
$[T]_{B_v'} = Q^{-1} [T]_{B_v} Q.$\\

Let $T : \mathbb{R}^2 \to \mathbb{R}^2$ be defined by $T\binom{a}{b} = \binom{3a - b}{a + 3b}$.
Let
\[
\mathcal{B}_V = \left\{ \binom{1}{1},\, \binom{-1}{1} \right\}, \qquad
\mathcal{B}_V' = \left\{ \binom{2}{4},\, \binom{3}{1} \right\}.
\]
Verify by direct computation that $[T]_{\mathcal{B}_V'} = Q^{-1}\,[T]_{\mathcal{B}_V}\, Q$, where $Q = [I_V]_{\mathcal{B}_V'}^{\mathcal{B}_V}$ is the change-of-coordinate matrix from $\mathcal{B}_V'$ to $\mathcal{B}_V$.

\noindent\\
\textbf{Solution.}\\
\quad\\
\begin{mdframed}[linewidth=1pt]
\textbf{Step 1: Compute $[T]_{\mathcal{B}_V}$.}

Apply $T$ to each basis vector:
\[
T\binom{1}{1} = \binom{2}{4}, \qquad T\binom{-1}{1} = \binom{-4}{2}.
\]

Express the images in $\mathcal{B}_V$-coordinates. Solve $\binom{2}{4} = a\binom{1}{1} + b\binom{-1}{1}$:
\[
a - b = 2, \quad a + b = 4 \implies a = 3,\; b = 1.
\]

Solve $\binom{-4}{2} = a\binom{1}{1} + b\binom{-1}{1}$:
\[
a - b = -4, \quad a + b = 2 \implies a = -1,\; b = 3.
\]

Therefore:
\[
[T]_{\mathcal{B}_V} = \begin{pmatrix} 3 & -1 \\ 1 & 3 \end{pmatrix}.
\]

\medskip
\textbf{Step 2: Compute $[T]_{\mathcal{B}_V'}$.}

Apply $T$ to each basis vector:
\[
T\binom{2}{4} = \binom{2}{14}, \qquad T\binom{3}{1} = \binom{8}{6}.
\]

Express in $\mathcal{B}_V'$-coordinates. Solve $\binom{2}{14} = a\binom{2}{4} + b\binom{3}{1}$:
\[
2a + 3b = 2, \quad 4a + b = 14.
\]
From the second equation, $b = 14 - 4a$. Substituting: $2a + 42 - 12a = 2$, so $a = 4$, $b = -2$.

Solve $\binom{8}{6} = a\binom{2}{4} + b\binom{3}{1}$:
\[
2a + 3b = 8, \quad 4a + b = 6.
\]
From the second equation, $b = 6 - 4a$. Substituting: $2a + 18 - 12a = 8$, so $a = 1$, $b = 2$.

Therefore:
\[
[T]_{\mathcal{B}_V'} = \begin{pmatrix} 4 & 1 \\ -2 & 2 \end{pmatrix}.
\]

\medskip
\textbf{Step 3: Compute $Q = [I_V]_{\mathcal{B}_V'}^{\mathcal{B}_V}$.}

Express each $\mathcal{B}_V'$ vector in $\mathcal{B}_V$-coordinates.

$\binom{2}{4} = a\binom{1}{1} + b\binom{-1}{1}$: $a - b = 2$, $a + b = 4$, so $a = 3$, $b = 1$.

$\binom{3}{1} = a\binom{1}{1} + b\binom{-1}{1}$: $a - b = 3$, $a + b = 1$, so $a = 2$, $b = -1$.

Therefore:
\[
Q = \begin{pmatrix} 3 & 2 \\ 1 & -1 \end{pmatrix}, \qquad
Q^{-1} = \frac{1}{\det Q}\begin{pmatrix} -1 & -2 \\ -1 & 3 \end{pmatrix} = \frac{1}{-5}\begin{pmatrix} -1 & -2 \\ -1 & 3 \end{pmatrix} = \frac{1}{5}\begin{pmatrix} 1 & 2 \\ 1 & -3 \end{pmatrix}.
\]

\medskip
\textbf{Step 4: Verify $Q^{-1}[T]_{\mathcal{B}_V}\, Q = [T]_{\mathcal{B}_V'}$.}

First compute $Q^{-1}[T]_{\mathcal{B}_V}$:
\[
\frac{1}{5}\begin{pmatrix} 1 & 2 \\ 1 & -3 \end{pmatrix}
\begin{pmatrix} 3 & -1 \\ 1 & 3 \end{pmatrix}
= \frac{1}{5}\begin{pmatrix} 5 & 5 \\ 0 & -10 \end{pmatrix}
= \begin{pmatrix} 1 & 1 \\ 0 & -2 \end{pmatrix}.
\]

Then multiply by $Q$:
\[
\begin{pmatrix} 1 & 1 \\ 0 & -2 \end{pmatrix}
\begin{pmatrix} 3 & 2 \\ 1 & -1 \end{pmatrix}
= \begin{pmatrix} 4 & 1 \\ -2 & 2 \end{pmatrix}
= [T]_{\mathcal{B}_V'}. \quad \checkmark
\]

\medskip
\textbf{Consistency check:} Both matrix representations have $\det = 10$ and $\operatorname{tr} = 6$, confirming they represent the same linear transformation. \checkmark
\end{mdframed}

\bigskip

%% ============================================================
%% Jordan canonical form problems
%% ============================================================

Find the Jordan canonical form of each of the following matrices.

\bigskip

%% ============================================================
%% Problem (a)
%% ============================================================

\needspace{4\baselineskip}
\hfill \break  $~\!\!$ \dotfill \medskip \\
\textbf{Question 2(a).}\\
Find the Jordan canonical form of
\[
A = \begin{pmatrix} 1 & 1 \\ -1 & 3 \end{pmatrix}.
\]

\noindent\\
\textbf{Solution.}\\
\quad\\
\begin{mdframed}[linewidth=1pt]
\textbf{Step 1: Characteristic polynomial.}

\begin{align*}
\det(A - \lambda I) &= \det\begin{pmatrix} 1-\lambda & 1 \\ -1 & 3-\lambda \end{pmatrix} \\
&= (1-\lambda)(3-\lambda) - (1)(-1) \\
&= \lambda^2 - 4\lambda + 3 + 1 \\
&= \lambda^2 - 4\lambda + 4 \\
&= (\lambda - 2)^2.
\end{align*}

So $\lambda = 2$ is the only eigenvalue, with algebraic multiplicity $2$.

\medskip
\textbf{Step 2: Geometric multiplicity.}

\[
A - 2I = \begin{pmatrix} -1 & 1 \\ -1 & 1 \end{pmatrix}.
\]

This matrix has rank $1$, so $\nu_1(2) = \dim(\ker(A - 2I)) = 2 - 1 = 1$.

\medskip
\textbf{Step 3: Jordan form.}

Since the algebraic multiplicity is $2$ and the geometric multiplicity is $1$, $A$ is \emph{not} diagonalizable. There must be a single Jordan block of size $2$:

\[
\boxed{J = \begin{pmatrix} 2 & 1 \\ 0 & 2 \end{pmatrix} = J_2(2).}
\]

We can verify using the theorem from lecture: $2\nu_1 - \nu_2 - \nu_0 = 2(1) - 2 - 0 = 0$ blocks of order $1$, and since the total size must be $2$, there is one block of order $2$.
\end{mdframed}

\bigskip

%% ============================================================
%% Problem (b)
%% ============================================================

\needspace{4\baselineskip}
\hfill \break  $~\!\!$ \dotfill \medskip \\
\textbf{Question 2(b).}\\
Find the Jordan canonical form of
\[
A = \begin{pmatrix} 1 & 2 \\ 3 & 2 \end{pmatrix}.
\]

\noindent\\
\textbf{Solution.}\\
\quad\\
\begin{mdframed}[linewidth=1pt]
\textbf{Step 1: Characteristic polynomial.}

\begin{align*}
\det(A - \lambda I) &= \det\begin{pmatrix} 1-\lambda & 2 \\ 3 & 2-\lambda \end{pmatrix} \\
&= (1-\lambda)(2-\lambda) - 6 \\
&= \lambda^2 - 3\lambda + 2 - 6 \\
&= \lambda^2 - 3\lambda - 4 \\
&= (\lambda - 4)(\lambda + 1).
\end{align*}

So the eigenvalues are $\lambda_1 = 4$ and $\lambda_2 = -1$, each with algebraic multiplicity $1$.

\medskip
\textbf{Step 2: Jordan form.}

Since both eigenvalues are distinct and have algebraic multiplicity $1$, the geometric multiplicity of each is also $1$. The matrix is diagonalizable, and its Jordan form consists of two $1 \times 1$ blocks:

\[
\boxed{J = \begin{pmatrix} 4 & 0 \\ 0 & -1 \end{pmatrix} = \operatorname{diag}\bigl(J_1(4),\, J_1(-1)\bigr).}
\]

\medskip
\textbf{Verification:} $\det J = (4)(-1) = -4 = \det A = (1)(2) - (2)(3) = -4$. $\operatorname{tr} J = 4 + (-1) = 3 = \operatorname{tr} A = 1 + 2 = 3$. \checkmark
\end{mdframed}

\bigskip

%% ============================================================
%% Problem (c)
%% ============================================================

\needspace{4\baselineskip}
\hfill \break  $~\!\!$ \dotfill \medskip \\
\textbf{Question 2(c).}\\
Find the Jordan canonical form of
\[
A = \begin{pmatrix} 11 & -4 & -5 \\ 21 & -8 & -11 \\ 3 & -1 & 0 \end{pmatrix}.
\]

\noindent\\
\textbf{Solution.}\\
\quad\\
\begin{mdframed}[linewidth=1pt]
\textbf{Step 1: Characteristic polynomial.}

We compute $\det(A - \lambda I)$. Expanding along the third row:

\begin{align*}
\det(A - \lambda I) &= \det\begin{pmatrix} 11-\lambda & -4 & -5 \\ 21 & -8-\lambda & -11 \\ 3 & -1 & -\lambda \end{pmatrix}.
\end{align*}

Expanding along the third row:
\begin{align*}
&= 3\det\begin{pmatrix} -4 & -5 \\ -8-\lambda & -11 \end{pmatrix} - (-1)\det\begin{pmatrix} 11-\lambda & -5 \\ 21 & -11 \end{pmatrix} + (-\lambda)\det\begin{pmatrix} 11-\lambda & -4 \\ 21 & -8-\lambda \end{pmatrix}.
\end{align*}

Computing each $2 \times 2$ determinant:
\begin{align*}
\det\begin{pmatrix} -4 & -5 \\ -8-\lambda & -11 \end{pmatrix} &= 44 - 5(8+\lambda) = 4 - 5\lambda, \\
\det\begin{pmatrix} 11-\lambda & -5 \\ 21 & -11 \end{pmatrix} &= -11(11-\lambda) + 105 = -16 + 11\lambda, \\
\det\begin{pmatrix} 11-\lambda & -4 \\ 21 & -8-\lambda \end{pmatrix} &= (11-\lambda)(-8-\lambda) + 84 = \lambda^2 - 3\lambda - 4.
\end{align*}

Combining:
\begin{align*}
\det(A - \lambda I) &= 3(4 - 5\lambda) + (-16 + 11\lambda) - \lambda(\lambda^2 - 3\lambda - 4) \\
&= 12 - 15\lambda - 16 + 11\lambda - \lambda^3 + 3\lambda^2 + 4\lambda \\
&= -\lambda^3 + 3\lambda^2 - 4.
\end{align*}

Testing $\lambda = -1$: $-(-1) + 3(1) - 4 = 1 + 3 - 4 = 0$. So $(\lambda + 1)$ is a factor.

Performing polynomial division:
\[
-\lambda^3 + 3\lambda^2 - 4 = -(\lambda + 1)(\lambda^2 - 4\lambda + 4) = -(\lambda + 1)(\lambda - 2)^2.
\]

Eigenvalues: $\lambda_1 = -1$ (algebraic multiplicity $1$), $\lambda_2 = 2$ (algebraic multiplicity $2$).

\medskip
\textbf{Step 2: Geometric multiplicity of $\lambda = 2$.}

\[
A - 2I = \begin{pmatrix} 9 & -4 & -5 \\ 21 & -10 & -11 \\ 3 & -1 & -2 \end{pmatrix}.
\]

Row reduce using $R_1 \leftarrow R_1 - 3R_3$ and $R_2 \leftarrow R_2 - 7R_3$:
\[
\begin{pmatrix} 0 & -1 & 1 \\ 0 & -3 & 3 \\ 3 & -1 & -2 \end{pmatrix}
\xrightarrow{R_2 \leftarrow R_2 - 3R_1}
\begin{pmatrix} 0 & -1 & 1 \\ 0 & 0 & 0 \\ 3 & -1 & -2 \end{pmatrix}.
\]

Rank $= 2$, so $\nu_1(2) = 3 - 2 = 1$.

\medskip
\textbf{Step 3: Jordan form.}

Since $\lambda = -1$ has algebraic multiplicity $1$, it contributes one $J_1(-1)$ block.

Since $\lambda = 2$ has algebraic multiplicity $2$ and geometric multiplicity $1$, it contributes one $J_2(2)$ block (not diagonalizable for this eigenvalue).

\[
\boxed{J = \begin{pmatrix} -1 & 0 & 0 \\ 0 & 2 & 1 \\ 0 & 0 & 2 \end{pmatrix} = \operatorname{diag}\bigl(J_1(-1),\, J_2(2)\bigr).}
\]

\medskip
\textbf{Verification:} $\det J = (-1)(2)(2) = -4$. $\det A = 11(-8 \cdot 0 - (-11)(-1)) - (-4)(21 \cdot 0 - (-11)(3)) + (-5)(21(-1) - (-8)(3)) = 11(-11) + 4(33) - 5(-21 + 24) = -121 + 132 - 15 = -4$. \checkmark

$\operatorname{tr} J = -1 + 2 + 2 = 3 = 11 + (-8) + 0 = \operatorname{tr} A$. \checkmark
\end{mdframed}

\bigskip

%% ============================================================
%% Problem (d)
%% ============================================================

\needspace{4\baselineskip}
\hfill \break  $~\!\!$ \dotfill \medskip \\
\textbf{Question 2(d).}\\
Find the Jordan canonical form of
\[
A = \begin{pmatrix} 2 & 1 & 0 & 0 \\ 0 & 2 & 1 & 0 \\ 0 & 0 & 3 & 0 \\ 0 & 1 & -1 & 3 \end{pmatrix}.
\]

\noindent\\
\textbf{Solution.}\\
\quad\\
\begin{mdframed}[linewidth=1pt]
\textbf{Step 1: Characteristic polynomial.}

\begin{align*}
A - \lambda I = \begin{pmatrix} 2-\lambda & 1 & 0 & 0 \\ 0 & 2-\lambda & 1 & 0 \\ 0 & 0 & 3-\lambda & 0 \\ 0 & 1 & -1 & 3-\lambda \end{pmatrix}.
\end{align*}

Expanding along the first column (only one nonzero entry in column 1):
\[
\det(A - \lambda I) = (2-\lambda) \det\begin{pmatrix} 2-\lambda & 1 & 0 \\ 0 & 3-\lambda & 0 \\ 1 & -1 & 3-\lambda \end{pmatrix}.
\]

Expanding the $3 \times 3$ determinant along the third column:
\[
\det\begin{pmatrix} 2-\lambda & 1 & 0 \\ 0 & 3-\lambda & 0 \\ 1 & -1 & 3-\lambda \end{pmatrix} = (3-\lambda)\det\begin{pmatrix} 2-\lambda & 1 \\ 0 & 3-\lambda \end{pmatrix} = (3-\lambda)(2-\lambda)(3-\lambda).
\]

Therefore:
\[
\det(A - \lambda I) = (2-\lambda)^2(3-\lambda)^2.
\]

Eigenvalues: $\lambda_1 = 2$ (algebraic multiplicity $2$), $\lambda_2 = 3$ (algebraic multiplicity $2$).

\medskip
\textbf{Step 2: Geometric multiplicity of $\lambda = 2$.}

\[
A - 2I = \begin{pmatrix} 0 & 1 & 0 & 0 \\ 0 & 0 & 1 & 0 \\ 0 & 0 & 1 & 0 \\ 0 & 1 & -1 & 1 \end{pmatrix}.
\]

Row reduce: $R_3 \leftarrow R_3 - R_2$, then $R_4 \leftarrow R_4 - R_1$, then $R_4 \leftarrow R_4 + R_2$:
\[
\begin{pmatrix} 0 & 1 & 0 & 0 \\ 0 & 0 & 1 & 0 \\ 0 & 0 & 0 & 0 \\ 0 & 0 & 0 & 1 \end{pmatrix}.
\]

Rank $= 3$, so $\nu_1(2) = 4 - 3 = 1$.

Since algebraic multiplicity $= 2$ and geometric multiplicity $= 1$, we get one $J_2(2)$ block.

\medskip
\textbf{Step 3: Geometric multiplicity of $\lambda = 3$.}

\[
A - 3I = \begin{pmatrix} -1 & 1 & 0 & 0 \\ 0 & -1 & 1 & 0 \\ 0 & 0 & 0 & 0 \\ 0 & 1 & -1 & 0 \end{pmatrix}.
\]

Row reduce: $R_4 \leftarrow R_4 + R_2$:
\[
\begin{pmatrix} -1 & 1 & 0 & 0 \\ 0 & -1 & 1 & 0 \\ 0 & 0 & 0 & 0 \\ 0 & 0 & 0 & 0 \end{pmatrix}.
\]

Rank $= 2$, so $\nu_1(3) = 4 - 2 = 2$.

Since algebraic multiplicity $= 2$ and geometric multiplicity $= 2$, $\lambda = 3$ is diagonalizable and contributes two $J_1(3)$ blocks.

\medskip
\textbf{Step 4: Jordan form.}

\[
\boxed{J = \begin{pmatrix} 2 & 1 & 0 & 0 \\ 0 & 2 & 0 & 0 \\ 0 & 0 & 3 & 0 \\ 0 & 0 & 0 & 3 \end{pmatrix} = \operatorname{diag}\bigl(J_2(2),\, J_1(3),\, J_1(3)\bigr).}
\]

\medskip
\textbf{Verification:}
\begin{itemize}
    \item $\det J = 2 \cdot 2 \cdot 3 \cdot 3 = 36$. $\det A = 2 \cdot 2 \cdot 3 \cdot 3 = 36$ (expanding along column 1, then noting the upper triangular-like structure). \checkmark
    \item $\operatorname{tr} J = 2 + 2 + 3 + 3 = 10 = 2 + 2 + 3 + 3 = \operatorname{tr} A$. \checkmark
    \item Characteristic polynomial: $(2-\lambda)^2(3-\lambda)^2$. \checkmark
\end{itemize}
\end{mdframed}

\end{document}
